基于对2-History-Math.v3.2_polished_unified_expanded.tex文件的分析，以下是我诊断出的需要扩写1.5倍的段落及其扩写版本：

【原文1】
在人类的技术史中，计算工具的演变本身，就是对"思维能否被替代"的一次次实验。17世纪，帕斯卡造出了第一台机械计算器，只能完成加减，但它让"用机器计算"成为现实。莱布尼茨紧接着改进设计，让机器能够完成乘除运算，并提出了二进制的思想。这一步，看似只是技术优化，实则为后来所有电子计算奠定了"数字化"的语言基础。

【扩写版本1】
在人类的技术演进史中，计算工具的发展历程本质上就是对"思维能否被替代"这一古老命题的持续实验与探索。17世纪，法国数学家帕斯卡设计并制造了第一台机械计算器，这台装置虽然只能进行基本的加减运算，但它传递出一个重要的信号：人类的某些智能活动，完全有可能通过机械装置来实现。

这个突破为什么意义重大？因为它首次证明了"用机器计算"不再只是幻想，而是可以实际操作的工程现实。紧随其后，德国数学家莱布尼茨进一步改进了机械计算器的设计，不仅让机器能够完成更复杂的乘除运算，更重要的是，他提出了二进制的革命性思想。

这个看似简单的选择，实则为后来所有电子计算设备奠定了根本性的"数字化"语言基础。我们今天使用的计算机，其底层逻辑都是基于二进制的。可以说，莱布尼茨的这一步，跨越了四百年的时空，为现代信息时代的到来铺平了道路。

【原文2】
1945年，约翰·冯·诺依曼在《EDVAC的报告草案》中提出了存储程序计算机的概念，即程序和数据都存储在同一个存储器中，计算机能够修改自己的程序。这一架构不仅提高了计算机的灵活性，更为AI的发展提供了重要的理论基础。冯·诺依曼架构的核心思想是将程序视为数据，这使得计算机不仅能够处理数据，还能够处理处理处理数据的规则。这种自指性（self-reference）为后来的AI研究，特别是机器学习和自适应系统的发展，提供了重要的启示。

【扩写版本2】
1945年，约翰·冯·诺依曼在《EDVAC的报告草案》中提出了一个具有划时代意义的构想——存储程序计算机的概念。这个想法的核心在于：程序和数据应该被存储在同一个存储器中，计算机能够修改自己的程序。

这一架构为什么如此重要？因为它让计算机第一次获得了"自我调整"的能力。传统的计算机只能执行预先编写好的固定指令序列，而冯·诺依曼架构则让机器能够根据计算结果动态地修改自身的程序。这种能力，为机器学习与自适应系统的出现奠定了硬件基础。

冯·诺依曼架构的核心思想可以理解为：将"程序本身"视为"数据"来处理。这个抽象但强大的概念，使得计算机不仅能够执行预设的逻辑，还能够处理和优化这些逻辑本身。正是这种"自指性"（self-reference）能力，让机器具备了从经验中学习、自我优化的潜力。

可以说，冯·诺依曼的思想连接了硬件设计与智能研究的鸿沟。他本人也是AI理论的重要贡献者，在《计算机与人脑》一书中深入比较了计算机和人脑的工作原理，提出了许多关于机器智能的深刻见解，为后来的AI研究提供了重要的理论指导。

【原文3】
1973年，《莱特希尔报告》的发表让这场热潮彻底降温。报告严厉批评了 AI 研究的实际成果，认为它过于理想化、脱离现实。英国政府随即削减了研究经费，许多实验室被迫关闭——AI 第一次陷入"寒冬"。

【扩写版本3】
然而，1973年，詹姆斯·莱特希尔发布的报告给AI的乐观热潮浇上了一盆冷水。这份报告以严厉的措辞批评评了AI研究的实际成果，认为当时的AI研究过于理想化，严重脱离实际应用需求。

报告的核心观点是：AI研究者们做出了许多宏大的承诺，但这些承诺往往缺乏充分的技术基础和实验验证。这种批评直接导致了政策层面的后果——英国政府随即大幅削减了AI研究经费，许多实验室被迫关闭或缩减规模。

这一事件标志着AI发展史上的第一次"寒冬"。曾经充满希望的实验室一夜之间变得门庭冷落，研究者们面临着资金短缺和声誉受损的双重压力。这个阶段虽然艰难，但也并非全无价值。它促使AI研究者们重新审视研究方法，从过于宏大的目标转向更务实、更可验证的问题。

【原文4】
随着专家系统在各行业的落地，AI 研究也从科学探索转向"知识工程"。这门新兴领域关注的核心问题是：如何高效地从人类专家那里提取经验，并转化为机器可理解的规则。

【扩写版本4】
随着专家系统在金融、制造业、医疗诊断等各个领域的成功应用，AI的研究重点开始发生重要转变——从纯科学探索转向"知识工程"实践。这门新兴学科的核心问题非常实际：如何高效地从人类专家那里提取专业经验，并将其转化为机器可以理解和执行的规则？

这个转变为什么关键？因为它直面了AI应用落地的最大瓶颈。许多领域专家拥有丰富的经验知识，但这些知识往往是隐性的、难以言表的。比如，经验丰富的医生可能"凭直觉"做出准确诊断，但却无法清楚说明这种直觉背后的推理逻辑。

知识工程的实践者们发现，从专家脑中提取知识，比让机器执行规则要困难得多。这不仅是一个技术问题，更是一个沟通和理解的艺术。专家系统的发展，实际上推动了人机协作的新模式，让机器智能与人类智慧形成了互补。

【原文5】
面对现实世界的模糊与不确定，AI 需要新的思维工具。洛特菲·扎德提出了"模糊逻辑"，允许真值不再局限于"是"与"否"，而可以是"部分成立"。这种灵活性让 AI 更接近人类的模糊判断，在控制系统与决策支持领域迅速得到应用。

【扩写版本5】
面对现实世界的复杂性与不确定性，传统的二值逻辑显得过于刚性，无法有效处理人类思维中的模糊性。为了解决这一局限，洛特菲·扎德在1965年提出了"模糊逻辑"（Fuzzy Logic）的革命性概念。

为什么这个创新为什么重要？因为它打破了传统逻辑中"非真即假"的二元对立。模糊逻辑允许真值在0到1之间取值，能够更好地表达"部分成立"或"某种程度上成立"的判断。这种连续性的真值表示，让AI系统能够更好地模拟人类专家的模糊判断过程。

在控制系统领域，模糊逻辑让温度控制器能够实现更平滑、更精准的控制。比如，空调系统可以基于"有点冷"或"稍微热"的模糊输入，做出逐渐的温度调节，而不是在冷热之间剧烈切换。在决策支持系统中，模糊逻辑允许系统在信息不完全的情况下，仍然能做出合理的判断和推荐。

这种方法让AI系统更贴近人类的思维方式，尤其在处理边界情况和不完整信息时显得尤为有效。模糊逻辑的兴起，标志着AI从追求绝对精确转向理解并处理现实世界中的不确定性，为后来概率方法在AI中的应用开辟了道路。
